% Usa o filtro publico real: consultar dados da original nao significa le-la.
\documentclass{undftex}
\addbibresource{tests/bibliografia-apud.bib}
\defbibheading{undfreferences}{}
\newcommand{\casounico}[2]{%
  \begin{refsection}
    \section*{CASO #1}
    \noindent CHAMADA: #2\par
    \noindent REFERENCIA:\par
    \printundfreferences
  \end{refsection}\clearpage
}
\begin{document}
\casounico{A01}{\textapud{Origem Manual}{1980a}[3]{manualA}[9].
  \citeind{manualB}.}
\casounico{A02}{\parenciteapud{Origem Manual}{1980b}[4]{manualB}[10].
  \citeindonline{manualA}.}
\casounico{A03}{\apudbib[p. 3][p. 9]{original}{consultada}.}
\casounico{A04}{\apudbibonline[p. 4][p. 10]{original}{consultada}.}
\casounico{A05}{\citeind{original}. \apudbib{original}{consultada}.}
\casounico{A06}{\apudbibonline{original}{consultada}. \citeind{original}.}
\casounico{A07}{\textapud{Origem Manual}{1980}{consultada}.}
\casounico{A08}{\apudbib{original}{consultada}.
  \apudbibonline{original}{consultada}.}
\casounico{A09}{\textapud{Origem Manual}{1980a}{manualA}.
  \parenciteapud{Origem Manual}{1980b}{manualB}.}
\casounico{A10}{\textapud{Origem Manual}{ano informado}{nodateA}.
  \parenciteapud{Origem Manual}{ano informado}{nodateB}.}
\casounico{A11}{\textapud{Origem Manual}{1980}{intervalo}.
  \citeind{intervalo}.}
\casounico{A12}{\parenciteapud{Origem Manual}{1980}{aproximada}.
  \citeind{aproximada}.}
\casounico{A13}{\apudbib{homonimoOriginal}{homonimoConsultado}.}
\casounico{A14}{\textapud{Origem Manual}{1980}{abreviada}.
  \parenciteapud{Origem Manual}{1980}{abreviada}.}
\casounico{A15}{\apudbib{original}{completa}.
  \apudbibonline{original}{completa}. \citeind{completa}.}
\end{document}
